\section{Descripción Funcional del Sistema}

\subsection{Secuencia de Operación}

Un ventilador centrífugo de impulsión toma aire exterior, lo hace atravesar una etapa de filtración de dos niveles (prefiltro MERV 8 + filtro final MERV 13-14) y lo inyecta directamente al laboratorio a 8 m/s (presión dinámica de inyección de 28 Pa en el sitio), sin red de ductos de distribución. El aporte continuo de 3 840 m$^3$/h (12 renovaciones/h) eleva la presión interior por encima de la exterior; el aire sale del recinto por tres rejillas de exfiltración de 353 mm $\times$ 336 mm provistas de malla anti-insectos y por las infiltraciones de la envolvente.

\subsection{Control de Presurización}

El set-point es +25 Pa con mínimo admisible de +12.5 Pa. El elemento final de control es un damper de alivio barométrico con contrapeso calibrable: permanece cerrado mientras el diferencial interior-exterior sea inferior al ajuste y abre progresivamente cuando lo supera, descargando el exceso de caudal y estabilizando la presión de sala. Dado que el contrapeso actúa sobre un diferencial de presión y no sobre caudal, el ajuste a +25 Pa es válido a cualquier altitud; la posición de equilibrio de apertura sí cambia con la densidad, por lo que la calibración final se ejecuta en comisionamiento con micromanómetro, práctica documentada por el fabricante \cite{halton_brd}. Con filtro limpio ($\Delta P$ total 95 Pa en el sitio) el ventilador entrega más presión de la necesaria y el damper trabaja más abierto; a medida que el filtro se carga hasta su $\Delta P$ final (190 Pa en el sitio), el damper cierra y mantiene la consigna. Esta compensación pasiva es la razón por la que el damper es obligatorio y no opcional (ver Sección 4.2).

\subsection{Lazo de Medición y Alarma}

Un transmisor de presión diferencial (4-20 mA, rango 0-62.5 Pa) mide el interior contra una referencia exterior protegida del viento y del ambiente clorado; sus relés (o el PLC/BMS del cliente) generan alarma baja a +12.5 Pa (pérdida de presurización: riesgo de ingreso de insectos, polvo y cloro) y alarma alta a +40 Pa (damper atascado), con retardo de 30-60 s para evitar falsas alarmas por apertura de puerta (dato típico de diseño). Un manómetro diferencial de lectura local (Magnehelic 0-60 Pa) en la puerta o pared del laboratorio permite verificación visual sin instrumentación auxiliar. El esquema corresponde al lazo documentado para presurización de edificios comerciales: sensor diferencial interior-exterior cuya señal supervisa el elemento de alivio \cite{trane_apn003}.

\subsection{Exclusión de Insectos y Polvo}

La barrera es triple: presión positiva permanente (flujo neto hacia afuera por cualquier abertura), filtración MERV 13-14 del aire de impulsión (E3 $\geq$ 90 \% para partículas de 3-10 \textmu m según ASHRAE 52.2 \cite{ashrae522_2017}) y malla anti-insectos inox 316 de 18$\times$18 (abertura $\approx$1 mm) en las rejillas de exfiltración, que impide el ingreso por las bocas de salida cuando el sistema está detenido.
